% titlepage.tex
% ============================================================
%  صفحه‌ی عنوان و شناسنامه‌ی کتاب
% ============================================================

\begin{titlepage}
\begin{tikzpicture}[remember picture, overlay]

  % --- Full-page background ---
  \fill[BgCream] (current page.south west) rectangle (current page.north east);

  % --- Top dark strip ---
  \fill[PrimaryDark]
    (current page.north west) rectangle
    ([yshift=-9cm]current page.north east);

  % --- Gold accent band ---
  \fill[AccentGold]
    ([yshift=-9cm]current page.north west) rectangle
    ([yshift=-9.4cm]current page.north east);

  % --- Decorative pattern on dark strip ---
  \foreach \i in {0,1,...,25}{
    \node[opacity=0.04, text=white, font=\fontsize{40}{40}\selectfont,
          rotate=30]
      at ([xshift=\i*1.2cm-3cm, yshift=-4cm]current page.north west)
      {◆};
  }
  \foreach \i in {0,1,...,25}{
    \node[opacity=0.03, text=white, font=\fontsize{30}{30}\selectfont,
          rotate=-15]
      at ([xshift=\i*1.5cm-2cm, yshift=-7cm]current page.north west)
      {◇};
  }

  % --- Main Title ---
  \node[anchor=east, text=white, font=\fontsize{30}{36}\selectfont\bfseries,
        text width=16cm, align=flush right]
    at ([yshift=-3.5cm, xshift=-2cm]current page.north east)
    {دولت، ملت و تنوع قومی};

  % --- Subtitle ---
  \node[anchor=east, text=AccentGold, font=\fontsize{16}{22}\selectfont,
        text width=16cm, align=flush right]
    at ([yshift=-5.8cm, xshift=-2cm]current page.north east)
    {مبانی نظری، الگوهای تطبیقی\\[4pt]
     و راهنمای اجرایی برای ایران};

  % --- English subtitle ---
  \node[anchor=east, text=white!80, font=\fontsize{11}{14}\selectfont\itshape]
    at ([yshift=-7.8cm, xshift=-2cm]current page.north east)
    {\en{State, Nation \& Ethnic Diversity: Theory, Comparative Models \& A Roadmap for Iran}};

  % --- Central decorative element ---
  \begin{scope}[shift={([yshift=-13cm]current page.north)}]
    \foreach \a in {0,45,...,315}{
      \draw[AccentGold, opacity=0.3, line width=0.5pt]
        (0,0) -- (\a:2.5);
    }
    \draw[AccentGold, line width=1.5pt] (0,0) circle (2.2);
    \draw[AccentGold, line width=0.6pt] (0,0) circle (2.5);
    \draw[AccentGold, line width=0.6pt] (0,0) circle (1.8);
    \node[text=PrimaryDark, font=\fontsize{9}{11}\selectfont, text width=3cm,
          align=center]
      at (0,0) {پژوهش تطبیقی\\سیاست تأسیسی\\مدیریت تنوع};
  \end{scope}

  % --- Bottom info area ---
  \node[anchor=south east, text=TextMid, font=\small,
        text width=8cm, align=flush right]
    at ([yshift=3cm, xshift=-2cm]current page.south east)
    {%
      \textcolor{PrimaryDark}{\bfseries تألیف:}\quad پژوهشگر سیاست تطبیقی\\[6pt]
      \textcolor{PrimaryDark}{\bfseries ناشر:}\quad مؤسسه‌ی مطالعات راهبردی\\[6pt]
      \textcolor{PrimaryDark}{\bfseries نوبت چاپ:}\quad ویرایش اول — ۱۴۰۴
    };

  % --- Bottom left decorative ---
  \fill[PrimaryDark, rounded corners=4pt]
    ([yshift=1.5cm, xshift=2cm]current page.south west)
    rectangle
    ([yshift=2cm, xshift=8cm]current page.south west);
  \node[text=AccentGold, font=\footnotesize\bfseries,
        anchor=west]
    at ([yshift=1.75cm, xshift=2.3cm]current page.south west)
    {\en{Comparative Politics \& Constitutional Design}};

  % --- Bottom strip ---
  \fill[PrimaryDark]
    (current page.south west) rectangle
    ([yshift=0.8cm]current page.south east);
  \fill[AccentGold]
    ([yshift=0.8cm]current page.south west) rectangle
    ([yshift=1.1cm]current page.south east);

\end{tikzpicture}
\end{titlepage}

% --- Blank back of title page ---
\clearpage
\thispagestyle{empty}
\begin{center}
\vspace*{\fill}
{\small\color{TextMid}
تمامی حقوق این اثر محفوظ است.\\[4pt]
بازنشر با ذکر منبع بلامانع است.\\[12pt]
\textcolor{PrimaryDark}{\bfseries شابک:}\quad
\en{978-600-XXX-XXX-X}\\[4pt]
\textcolor{PrimaryDark}{\bfseries نوبت چاپ:}\quad اول — بهار ۱۴۰۴
}
\vspace*{\fill}
\end{center}
\clearpage